% ============================================================
% Section 147: n⁺n高-低结的欧姆导电性能
% ============================================================
\section{半导体物理：n⁺n高-低结的欧姆导电性能}
\label{sec:hl-junction-ohmic}

\begin{modelbox}[题目：n⁺n高-低结的欧姆导电特性]
在半导体器件制造中，常遇到低掺杂半导体引线问题，一般采用金属--n⁺型半导体--n型半导体的结构来解决，其中n⁺n称高-低结（H-L junction）。试分别画出给n⁺n结施加极性相反电压时的能带图，从而说明其欧姆导电性能。
\end{modelbox}

\begin{tipbox}[考点快速分析]
\begin{itemize}
    \item \textbf{热平衡}：电子从费米能级高的n⁺区流向n区，界面形成空间电荷区
    \item \textbf{n⁺区}：极薄的正空间电荷层（掺杂浓度高）
    \item \textbf{n区}：较宽的负空间电荷层（电子积累，形成势阱）
    \item \textbf{外加电压}：主要降落在高阻的轻掺杂n区，H-L结势垒变化极小
\end{itemize}
\end{tipbox}

\begin{derivationbox}[标准解题过程]
\textbf{1. 热平衡能带图}

n⁺区为重掺杂，费米能级 $E_{F,n^+}$ 靠近导带底；n区为轻掺杂，$E_{F,n}$ 相对较低。接触后电子从n⁺区流向n区：
\begin{itemize}
    \item n⁺区一侧：电子流失，留下电离施主正电荷，形成极薄的正空间电荷层。
    \item n区一侧：电子积累，形成负空间电荷层（势阱）。
\end{itemize}
能带在空间电荷区弯曲，内建电势 $eV_D=E_{F,n^+}-E_{F,n}$ 较小。空间电荷区主要落在轻掺杂n区一侧。

\textbf{2. 外加正向偏压（n区接正，n⁺区接负）}

外加电场与内建电场方向相反，H-L结势垒高度略微降低。电压主要降落在n区，n区能带向下倾斜。n⁺区的电子只需克服极小的势垒即可进入n区的高电导势阱区，然后在外加电场推动下漂移通过n区到达正电极。

\textbf{3. 外加反向偏压（n⁺区接正，n区接负）}

外加电场与内建电场方向相同，H-L结势垒高度略微升高。电压仍主要降落在n区，n区能带向上倾斜。n区的电子在外加电场推动下漂移通过n区，进入高电导势阱区，然后毫无阻碍地滑入n⁺区到达正电极。

\textbf{4. 结论}

无论外加电压极性如何，电流均主要由轻掺杂n区的体电阻限制，伏安特性为\textbf{线性且对称}，呈现欧姆导电行为。H-L结本身不引入整流效应。

\textbf{5. 金属-n⁺接触的欧姆特性（隧道效应）}

金属与重掺杂半导体接触时，耗尽层极薄（$x_D\propto1/\sqrt{N_D}$，$N_D\sim10^{19}\,\text{cm}^{-3}$ 时势垒宽度可薄至 $100\,\text{Å}$ 以下）。载流子可量子隧穿通过势垒，隧穿电流对偏压极性不敏感，正反向电流均很大且对称，形成良好的\textbf{欧姆接触}。

因此，金属--n⁺--n结构完美地解决了低掺杂半导体引线的欧姆接触问题。
\end{derivationbox}

\begin{formulabox}[答案汇总]
\begin{center}
\begin{tabular}{lll}
\toprule
\textbf{偏置情况} & \textbf{能带变化} & \textbf{电流机制} \\
\midrule
热平衡 & n⁺侧微小势垒，n侧势阱 & 无净电流 \\
n区接正 & n区能带下倾，势垒略降 & 电子越过小势垒，经n区漂移 \\
n⁺区接正 & n区能带上倾，势垒略升 & 电子漂移进入势阱，滑入n⁺区 \\
\bottomrule
\end{tabular}
\end{center}
\end{formulabox}

\begin{insightbox}[物理图像]
\begin{itemize}
    \item 高-低结本身不是整流结，其能带弯曲形成的是对多子而言的微小势垒或势阱，正反向均不阻碍多子流动。
    \item 金属与重掺杂半导体接触利用\textbf{隧道效应}实现欧姆接触，是半导体器件电极制备的标准工艺。
    \item ``金属--n⁺--n''结构将欧姆接触的难题从低掺杂区转移到了重掺杂区，既保证了良好的欧姆特性，又不对低掺杂有源区引入额外势垒。
\end{itemize}
\end{insightbox}
